\section{Dataset Description and Splitting}
\label{ch:dataset}

\subsection{Source and selected splits}
MS COCO is a dataset of everyday scenes with object-instance annotations and human-written captions \citep{lin2014coco,cocodocs}. The project uses the 2017 training and validation sources. The full train2017 annotation audit contains 118,287 images, 860,001 object instances, and 591,753 captions. The val2017 audit contains 5,000 images, 36,781 object instances, and 25,014 captions.

The frozen project subset contains 10,000 train images selected from train2017, and 100 validation plus 1,000 test images selected from val2017. Every selected image has exactly five reference captions. The train split supplies the retrieval knowledge base; validation is available for tuning, and test is held out for evaluation.

\begin{table}[htbp]
\centering
\caption{Frozen project splits and reference-caption counts.}
\label{tab:dataset-splits}
\rowcolors{2}{ProjectGrey}{white}
\begin{tabular}{@{}lrrr@{}}
\toprule
\textbf{Split} & \textbf{Images} & \textbf{Captions} & \textbf{Refs./image}\\
\midrule
Train & 10,000 & 50,000 & 5\\
Validation & 100 & 500 & 5\\
Test & 1,000 & 5,000 & 5\\
\midrule
\textbf{Total} & \textbf{11,100} & \textbf{55,500} & \textbf{5}\\
\bottomrule
\end{tabular}
\end{table}


The manifests are frozen in JSON and CSV under \code{data/selected/}. A rarity-first multilabel selection procedure uses seed 20261106, with deterministic offsets for the train and validation selections. Selection favors coverage of object categories, annotation-count complexity, and object-count groupings; it does not guarantee representativeness.

\subsection{References and leakage controls}
Each selected image has five references. Caption records preserve raw and cleaned fields; the caption index uses raw training captions only. Both FAISS indexes contain training items. The manifests have no image-ID overlap, and retrieval removes a query's own ID. Test references must remain outside retrieval and generation evidence.

\subsection{Test-size note}
The original assignment requirements specify 1,000 test images in the detailed dataset and preprocessing instructions. Although one occurrence contains the typographical value \code{1,00}, the detailed requirements and data-audit instructions establish 1,000 as the intended test-set size. Therefore, this project uses 1,000 test images, consistent with the detailed requirements.
